% Conclusion
\section{Conclusion}
\label{sec:conclusion}

We proposed \isr (\ISR), a configuration-space locality-weighted sector framework for
regulating toy eternal-inflation ensembles.
In a stochastic landscape simulation with 10,000 trajectories and six
vacuum basins, \ISR shifts probability mass toward the basin local to the
observable patch and improves trajectory-cutoff stability relative to a naive
global census by approximately a factor of two in KL divergence.
Across 100 sigma-by-seed stress tests, \ISR outperforms global census in
98 cases.
A shifted-well landscape perturbation suggests that the improvement is not
solely an artifact of the original well placement.

An ell-sensitivity scan (15 values from 0.05 to 10.0) reveals a sweet spot at
$\ell \approx 0.5$--$1.5$ where $N_{\rm eff}$ is large enough for statistical
power and basin-exception rates remain below 30\%.
A randomized-distance null control shows that the improvement is not a generic
artifact of kernel smoothing; it depends on preserving the field-space distance
structure encoded by the toy landscape.
In the non-degenerate region of the sigma-by-ell grid, ISR improves over global
census across every tested cell. Degenerate small-ell cells are excluded from
this evidence claim and reported separately in Appendix~\ref{app:degenerate}.

These results do not solve the eternal-inflation measure problem.
They demonstrate a reproducible toy framework for testing locality-weighted
measures, with built-in diagnostics that expose when and why the weighting
succeeds or fails.
The code, data, and paper source are available at
\url{https://github.com/jenholm/physics_models/tree/main/inflationary-sector-renormalization}.
